% ============================================================
% BME VIK TDK abstract — VARIANT 1: problem-first, method-as-solution
% (GNN + Adapter + LoRA + Contrastive loss)
%
% This is one of two abstract framings for the supervisors to compare
% and choose between. This variant leads with the missing-abstract
% problem and presents the architecture as the proposed solution.
% The second variant (falsifiability/grounding-question framing, per
% tdk-winning-strategy.md) is a separate file.
%
% Structure follows the official one-page sample:
% https://portal.vik.bme.hu/files/00000047.pdf
%   - max 1 page, min ~20 lines of body text
%   - NO header, footer, or page numbers
%   - figures/equations only if strictly necessary
%   - own contribution must be clearly separated from cited work
%   - body must cover: methodology, key results, conclusions
%
% NOTE ON LANGUAGE: the official TDK conference language is Hungarian
% (English talks require prior agreement with the section chair, per
% the official sample). This English draft is for internal review with
% Bálint/Dani; translate to Hungarian before submission unless an
% English exception is confirmed.
% ============================================================

\documentclass[12pt]{article}

\usepackage[utf8]{inputenc}
\usepackage[T1]{fontenc}
\usepackage[english]{babel}
\usepackage[a4paper, margin=2.5cm]{geometry}
\usepackage{mathptmx}      % Times-like font
\usepackage{hyperref}
\usepackage{enumitem}

\pagestyle{empty}          % <- enforces "no header/footer/page numbers"
\setlength{\parindent}{0pt}
\setlength{\parskip}{0.6em}

\begin{document}

% ---------- Title ----------
\begin{center}
    {\Large\bfseries Graph Conditioned Abstract Generation for Cross-Field Synthesis}
\end{center}

\vspace{0.8em}

% ---------- Authors ----------
\begin{center}
    Gema Csanád IV. Inf, gemacsanad@gmail.com \\
    
\end{center}

\vspace{0.4em}

% ---------- Consultant ----------
\begin{center}
    Consultant: Unyi Dániel 
\end{center}

\vspace{1.2em}

% ==========================================================
% BODY
% ==========================================================

A citation graph represents scientific literature as a network: each
node is a publication, each edge is a citation from one paper to
another, and every node carries that paper's abstract as its text
content. We study a specific task on this graph: given a target
paper's reference list, but withholding the target paper's own abstract
entirely, generate what that abstract says. The model never sees the
target paper's own words, it only sees which other papers it points
to, and what those papers say.

This is a meaningful test of whether a paper's contribution is
legible purely from its position in the graph. A paper's abstract
typically connects and synthesises ideas from several otherwise
unrelated fields, and its reference list is the only trace of
that synthesis available before the abstract itself exists.
Reconstructing the abstract from references alone tests whether this
cross-field connection is recoverable from the graph, independently
of the target paper's own text.

We approach this task with a system that combines a graph neural
network (GNN) with a language model. The language model generates the
target abstract conditioned on two things: the plain text of the
directly-cited papers' abstracts, and a structural embedding produced
by the GNN. The GNN aggregates information from beyond the
directly-cited papers themselves, capturing how these papers connect
to each other and to the wider citation network. An adapter
module projects this structural embedding into soft prefix tokens,
injected into every attention layer of the language model, which is
fine-tuned with LoRA. A contrastive alignment loss between the GNN
embedding and a frozen text encoding of the true target abstract
encourages the structural embedding to be directly predictive of the
target's content during training.

We construct the task on ogbn-arxiv-TA, a citation graph of computer
science papers extended with cross-domain reference nodes, so that
the cross-field citations central to the task are actually present
in a paper's neighbourhood. Generation quality is measured with
ROUGE,  BERTScore, and an LLM-as-judge protocol, comparing
the full model against a text-only baseline that
receives only the same directly-cited papers' abstracts. This
isolates whether the structural embedding adds predictive signal
beyond what is already visible in the directly-cited papers' own
text.

\vspace{1em}


\end{document}
